\section{Scope Boundaries, Non-Claims \& Future Work (ST-017)}
\label{sec:limitations}

\subsection{Non-Claims \& Scope Boundaries}
\begin{enumerate}
    \item \textbf{Model Scope:} ST-016 establishes an algebraic framework and
    witness-based demonstration that $C_{\text{contract}}$, $C_{\text{uncertainty}}$,
    and $C_{\text{temporal}}$ are necessary in formalized finite-instance scenarios.
    It does \emph{not} claim universal necessity for all possible governed runtimes,
    continuous control systems, or runtime models other than the ST-016 formal model.

    \item \textbf{Capability Semantics:} The formal Lean model treats
    $C_{\text{contract}}$, $C_{\text{uncertainty}}$, and $C_{\text{temporal}}$
    as abstract labels. The semantic connection to contract soundness, uncertainty
    bounding, and temporal consistency is established through finite-instance
    examples. General semantic grounding is reserved for future work.

    \item \textbf{Axiom0:} The structural sufficiency module postulates
    $\text{kernel}_D(x,y) \leftrightarrow K_D(x,y)$. This is a modelling
    assumption characterizing the class of decision problems studied, not a
    derived theorem.

    \item \textbf{Witness Bridge Automation:} The formal witness validation schema
    requires manual construction of \texttt{WitnessArtifact} records from
    observed TypeScript traces. No automatic extraction mechanism from runtime
    executions to Lean proofs exists in ST-016.

    \item \textbf{Transportability (ST-017 Boundary):} Transporting certified
    witness artifacts across heterogeneous runtimes ($M_1 \sim M_2$) is reserved
    for \textbf{ST-017}.
\end{enumerate}

\subsection{Open Research Direction: ST016\_Conjecture}
The following claim is \textbf{not proven} in this work and constitutes the
central open research direction for ST-017:

\begin{quote}
\textit{ST016 Conjecture:} For any runtime $M$ with $\mathcal{C} =
\{C_{\text{contract}}, C_{\text{uncertainty}}, C_{\text{temporal}}\}$
and any optimal policy $\pi^*$, if $M$ is decision-preserving then $M$ is
minimal (i.e., no capability can be removed while preserving $\pi^*$).
\end{quote}

Proving this conjecture would require establishing formal semantic axioms for
each capability class and demonstrating irreducibility in a model-independent
way, which is the core research question of ST-017.

\subsection{Conclusion}
ST-016 establishes an algebraic Lean 4 framework for reasoning about minimal
decision-preserving runtime governance, backed by formalized witness scenarios
and an audited replication package. The central open conjecture and semantic
grounding of capability classes provide the foundation for the ST-017 research
line on transportability and generalization.
